\documentclass[11pt, a4paper]{article}
\usepackage[T1]{fontenc}
\usepackage[utf8]{inputenc}
\usepackage[a4paper, top=2.5cm, bottom=2.5cm, left=2.5cm, right=2.5cm]{geometry}
\usepackage{booktabs}
\usepackage{tabularx}
\usepackage{enumitem}
\usepackage[hidelinks]{hyperref}
\usepackage{parskip}
\usepackage{graphicx}
\usepackage{amsmath}
\usepackage{float}
\setlength{\parindent}{0pt}
\pagestyle{plain}

\begin{document}

% ------------------------------------------------------------------------------
% PORTADA
% ------------------------------------------------------------------------------
\begin{center}
{\Large \textbf{Pontificia Universidad Javeriana}}\\[0.3cm]
Facultad de Ciencias Economicas y Administrativas\\[0.3cm]
Estadistica -- 2026-II\\[1cm]
{\LARGE \textbf{Genero y carreras de economia: ?`existe una brecha de genero en competencias cuantitativas?}}\\[0.5cm]
{\large Informe Final del Proyecto}\\[1cm]
\textbf{Grupo 4}\\[0.3cm]
Carrisosa Tovar Jhon Neyder\\
Hernandez Avila Luna Sofia\\
Gutierrez Lopez Yaddy Valentina\\
Martinez Marentes Juliana\\
Perez Abril Sebastian Andres\\[1cm]
27 de noviembre de 2026
\end{center}

\newpage
\tableofcontents
\newpage

% ------------------------------------------------------------------------------
\section{Introduccion}
% ------------------------------------------------------------------------------
% Contexto: brecha de genero en educacion superior, particularmente en
% competencias cuantitativas y de ingles en carreras economicas.
% Pregunta: la brecha de genero en razonamiento cuantitativo depende del
% nivel de habilidad previa (SB11)? Es diferente en economia vs ingenieria?
% Relevancia social: equidad de genero en el mercado laboral de economistas.

% ------------------------------------------------------------------------------
\section{Datos y Metodologia}
% ------------------------------------------------------------------------------
% Fuente: cruce Saber 11 -- Saber Pro (ICFES).
% Archivos: cruce\_saber11\_saberpro.parquet, saber11.parquet, saber\_pro.parquet.
% Filtro por carrera: ESTU\_PRGM\_ACADEMICO matching ADMINISTRA|ECONOM|CONTAD|NEGOCIOS.
% Variable dependiente 1: MOD\_RAZONA\_CUANTITAT\_PUNT.
% Variable dependiente 2: MOD\_INGLES\_PUNT.
% Predictores: PUNT\_GLOBAL\_sb11, ESTU\_GENERO\_sb11, FAMI\_ESTRATOVIVIENDA\_sb11,
%   INST\_ORIGEN\_sbpro.
% Modelo 1 (aditivo econ), Modelo 2 (interaccion genero * SB11, econ),
%   Modelo 3 (ingles como Y, econ), Modelo 4 (interaccion en ingenieria).
% Justificacion de la comparacion entre carreras economicas e ingenieria.

% ------------------------------------------------------------------------------
\section{Analisis Exploratorio (Fase 1)}
% ------------------------------------------------------------------------------
% Distribucion de razonamiento cuantitativo e ingles por genero.
% Scatter SB11 vs razonamiento cuantitativo coloreado por genero.
% Composicion de genero en carreras economicas vs ingenieria.

% ------------------------------------------------------------------------------
\section{Intervalos de Confianza y Pruebas de Hipotesis (Fase 2)}
% ------------------------------------------------------------------------------
% IC para la media de razonamiento cuantitativo por genero en carreras econ.
% Prueba t para diferencia de medias entre hombres y mujeres.
% Prueba de Wilcoxon como alternativa no parametrica.
% Prueba chi-cuadrado para composicion de genero por tipo de carrera.

% ------------------------------------------------------------------------------
\section{Modelo de Regresion (Fase 3)}
% ------------------------------------------------------------------------------
% Modelo 1: razon. cuantitativo $\sim$ SB11 + genero + estrato + inst\_origen (econ).
% Modelo 2: razon. cuantitativo $\sim$ SB11 * genero + estrato + inst\_origen (econ).
% Modelo 3: ingles $\sim$ SB11 * genero + estrato + inst\_origen (econ).
% Modelo 4: razon. cuantitativo $\sim$ SB11 * genero + estrato + inst\_origen (ing).
% Tabla de coeficientes con IC al 95\%.
% Comparacion de modelos: R2 ajustado, AIC, BIC.
% Predicciones: mujer vs hombre al percentil 50 de SB11.

% ------------------------------------------------------------------------------
\section{Diagnosticos y Robustez}
% ------------------------------------------------------------------------------
% Diagnosticos para cada modelo: residuales vs ajustados, QQ-plot.
% Test de Breusch-Pagan (homocedasticidad).
% VIF (multicolinealidad, incluyendo manejo de interacciones).
% Distancia de Cook (observaciones influyentes).
% Errores robustos HC1 para los 4 modelos.

% ------------------------------------------------------------------------------
\section{Resultados y Discusion}
% ------------------------------------------------------------------------------
% La interaccion genero * SB11 es significativa?
% Comparacion de la brecha entre carreras economicas e ingenieria.
% Graficos de interaccion, coeficientes, predicciones.
% Reflexion sobre equidad de genero en formacion economica y cuantitativa.

% ------------------------------------------------------------------------------
\section{Conclusiones}
% ------------------------------------------------------------------------------
% Respuesta a la pregunta de investigacion.
% Limitaciones del estudio.
% Implicaciones para politica de genero en educacion superior.

% ------------------------------------------------------------------------------
\section*{Referencias}
% ------------------------------------------------------------------------------
% ICFES, literatura sobre brecha de genero en STEM y ciencias economicas.

% ------------------------------------------------------------------------------
\section*{Anexo: Codigo R}
% ------------------------------------------------------------------------------
% Incluir codigo R utilizado (fase3\_grupo4.R) o referenciarlo.

\end{document}
